\documentclass[a4paper]{article}

\usepackage[ngerman]{babel}
\usepackage[T1]{fontenc}
\usepackage[utf8]{inputenc}
\usepackage{graphicx}
\usepackage{amsmath,amssymb}
\usepackage{hyperref}

\title{}
\date{}

\begin{document}
    
    \maketitle
    
    
    \section{Die Rolle von \textit{doch}}
        
        \subsection{\textit{doch} als Modalpartikel (Adverb)}
            
            \begin{itemize}
                \item \textit{doch} ist meistens eine Modalpartikel.
                \item Eine Modalpartikel zeigt die Haltung oder Erwartung des Sprechers.
                \item Sie ändert nicht die Grammatik des Satzes, sondern den Ton.
                \item Man kann sie oft weglassen, aber die Bedeutung im Sinne der Haltung verändert sich.
            \end{itemize}
            
            \textbf{Beispiel:}
            
            Du kommst doch morgen.
            
            Hier zeigt \textit{doch}, dass der Sprecher erwartet, dass die Person morgen kommt.
            
            \bigskip
        
        \subsection{\textit{doch} als Antwortpartikel}
            
            \begin{itemize}
                \item \textit{doch} kann allein stehen.
                \item Es wird benutzt, um einer negativen Aussage zu widersprechen.
            \end{itemize}
            
            \textbf{Beispiel:}
            
            A: Du hast kein Geld.
            
            B: Doch!
            
            Hier bedeutet \textit{doch}: Die Aussage ist falsch.
            
            \bigskip
        
        \subsection{Zusammenfassung}
            
            \begin{itemize}
                \item Wortart: Adverb
                \item Meistens Funktion: Modalpartikel
                \item Manchmal Funktion: Antwortpartikel
            \end{itemize}

\end{document}
